\subsection{Lesion-level specificity}
\label{sec:res_lesion}

A metric intended for radiology should respond more strongly to changes in
pathological tissue than to equivalent changes elsewhere. We test this on the
LGG collection \citep{buda2019association}, which provides expert tumour masks.
Gaussian noise of scale $\sigma$ is applied either inside the annotated tumour
region or inside an area-matched control region, and we report the ratio of the
induced discrepancy in the two conditions. Two controls are used: a mirrored
region at the contralateral position, and a texture-matched region selected to
have comparable gradient energy ($4.97\times10^{4}$ for tumour,
$4.77\times10^{4}$ mirrored, $5.33\times10^{4}$ texture-matched), which
controls for the possibility that the metric merely responds to local
structural complexity.

The RadioDINO MMD is lesion-selective where an Inception-based MMD is not, but
only in the upper part of the sweep (Table~\ref{tab:lesion}). At $\sigma = 60$
and $\sigma = 90$ the bootstrap interval excludes unity ($1.57$,
$[1.03, 1.97]$ and $1.39$, $[1.06, 1.71]$ against the mirrored control); at
$\sigma = 45$ the point estimate is larger ($1.71$) but the interval
$[0.94, 2.14]$ includes unity. At $\sigma \le 30$ the control-region response is
near zero and the ratio is numerically degenerate, so we report no value. The
Inception baseline stays within $1.01$--$1.08$ throughout, with intervals
covering unity at every $\sigma$.

We state this as suggestive rather than established, for two reasons. The
effect is individually significant at only two of the six perturbation levels
tested. More importantly, the perturbation is additive Gaussian noise confined
to a mask, which is a proxy for pathology-relevant change rather than an
instance of it: it establishes that the embedding weights tumour regions
differently from matched healthy regions, but not that it tracks clinically
meaningful pathology. Establishing the latter would require
pathology-preserving edits, lesion insertion or removal, or agreement with
radiologist judgement, none of which we attempt here.

\begin{table}[t]
\centering
\caption{Lesion selectivity: ratio of the discrepancy induced by perturbing
tumour tissue to that induced by perturbing an area-matched control region.
A value of $1.0$ means no lesion selectivity. Bootstrap $95\%$ intervals;
$\dagger$ marks intervals excluding unity. Values at $\sigma \le 30$ are
omitted because the control response is near zero and the ratio is degenerate.}
\label{tab:lesion}
\small
\begin{tabular}{rcccc}
\toprule
& \multicolumn{2}{c}{RadioDINO MMD} & \multicolumn{2}{c}{Inception MMD} \\
\cmidrule(lr){2-3}\cmidrule(lr){4-5}
$\sigma$ & mirrored & texture-matched & mirrored & texture-matched \\
\midrule
$45$ & $1.71\,[0.94, 2.14]$ & $1.81\,[0.94, 2.12]$
     & $1.08\,[0.92, 1.24]$ & $1.03\,[0.87, 1.16]$ \\
$60$ & $1.57\,[1.03, 1.97]^{\dagger}$ & $1.59\,[1.02, 1.90]^{\dagger}$
     & $1.08\,[0.93, 1.22]$ & $1.05\,[0.91, 1.17]$ \\
$90$ & $1.39\,[1.06, 1.71]^{\dagger}$ & $1.40\,[1.03, 1.69]^{\dagger}$
     & $1.06\,[0.93, 1.19]$ & $1.06\,[0.94, 1.17]$ \\
\bottomrule
\end{tabular}
\end{table}

\begin{figure}[t]
\centering
\includegraphics[width=0.9\linewidth]{Images/lgg_noise_sweep.png}
\caption{Lesion selectivity across the perturbation sweep on LGG, with
area-matched mirrored and texture-matched controls. Shaded bands are bootstrap
$95\%$ intervals. The shaded region at low $\sigma$ marks where the
control-region response is near zero and the ratio is not measurable.}
\label{fig:lesion}
\end{figure}
